\chapter{Designs}

\section{System Architecture Overview}
ScamGuard follows the usual three-tier layout. At the top, the \textbf{presentation tier} is a Next.js front end, and every API call goes out through its own origin. In the middle, the \textbf{application tier} is a FastAPI back end on Render; one process runs the detector, the intent rules, the cross-domain models, sign-in and the safety tools. At the bottom, the \textbf{data tier} is a PostgreSQL database on Neon holding users, scans, cases and reports. Training happens offline. The trained models are saved as versioned files, which the back end loads a single time when it starts. The deployed application can be tried at \url{https://scamguard-frontend-r6t7.onrender.com}.

The system therefore has an offline part and an online part, shown in Figure~\ref{fig:arch}. The shaded blocks run offline: the cross-domain notebook, the training on ScamGuard-IN-2026 and the model registry. Training writes versioned model files to the registry, and the back end loads them a single time when it starts. The online path follows each request through the web app, the FastAPI back end, the cleaner and tokenizer, the TF-IDF and logistic-regression classifier, the intent rules, the fusion step and the explanation layer, and ends by saving the scan in PostgreSQL. The AI copilot and the known-pattern retrieval (52 entries) feed the explanation layer, and the cross-domain page calls the two stored model pairs directly.

\begin{figure}[H]
\centering
\resizebox{\linewidth}{!}{\begin{tikzpicture}[
  box/.style={draw, rounded corners=2pt, align=center, font=\small, minimum width=3.8cm, minimum height=12mm, inner sep=4pt, fill=white},
  off/.style={box, fill=gray!15},
  arr/.style={-{Stealth[length=2mm]}, thick}]
\def\dx{4.6}
% offline row
\node[off] (xnb)   at (0,0)        {Cross-domain notebook\\(UCI + India)};
\node[box] (xdm)   at (\dx,0)      {Cross-domain models\\(baseline, adapted)};
\node[off] (train) at (2*\dx,0)    {Offline training on\\ScamGuard-IN-2026};
\node[off] (reg)   at (3*\dx,0)    {Model registry\\(versioned files)};
% request row
\node[box] (ui)    at (0,-2.2)     {Web app\\(Next.js)};
\node[box] (api)   at (\dx,-2.2)   {FastAPI\\back end};
\node[box] (pre)   at (2*\dx,-2.2) {Cleaner +\\tokenizer};
\node[box] (clf)   at (3*\dx,-2.2) {TF-IDF +\\logistic regression};
% response row
\node[box] (db)    at (0,-4.4)     {PostgreSQL\\database};
\node[box] (expl)  at (\dx,-4.4)   {Scam X-Ray,\\TrustLens};
\node[box] (fuse)  at (2*\dx,-4.4) {Fusion, guards,\\risk score};
\node[box] (rules) at (3*\dx,-4.4) {Intent and\\protective rules};
% helpers
\node[box] (cop)   at (\dx,-6.6)   {AI copilot\\(Gemini, capped)};
\node[box] (rag)   at (2*\dx,-6.6) {Known-pattern\\retrieval (52 entries)};
\draw[arr] (xnb) -- (xdm); \draw[arr] (train) -- (reg); \draw[arr] (reg) -- (clf);
\draw[arr, <->] (api) -- (xdm);
\draw[arr] (ui) -- (api); \draw[arr] (api) -- (pre); \draw[arr] (pre) -- (clf);
\draw[arr] (clf) -- (rules); \draw[arr] (rules) -- (fuse); \draw[arr] (fuse) -- (expl);
\draw[arr] (expl) -- (db); \draw[arr] (db) -- (ui);
\draw[arr] (rag) -- (expl); \draw[arr] (expl) -- (cop);
\end{tikzpicture}
}
\caption{Architecture of the deployed system. Shaded parts run offline.}
\label{fig:arch}
\end{figure}

\section{Module/Component Design}
Table~\ref{tab:subsys} breaks ScamGuard into modules and gives each one's job and the way other parts talk to it.

\begin{table}[H]
\centering
\caption{Modules of ScamGuard and their interfaces}
\label{tab:subsys}
\begin{tabular}{|L{3.2cm}|L{5.5cm}|L{5.2cm}|}
\hline
\textbf{Module} & \textbf{Responsibility} & \textbf{Interface} \\
\hline
Dataset builder & Gathers the sources, audits labels, groups campaigns and splits the data & Python script that writes one CSV file \\
\hline
Training pipeline & Trains and compares models, then registers the chosen one & Model registry (vectorizer and model files) \\
\hline
Detector & Cleans, classifies, applies rules, fuses scores and explains & Python function \texttt{predict(text)} \\
\hline
Cross-domain study & Audits and cleans the UCI and Indian data, runs the seven experiments over three seeds, saves results and models & Jupyter notebook producing CSV tables, figures, \texttt{results.json} and joblib models \\
\hline
Cross-domain service & Serves the baseline and adapted models, with explanations & \texttt{POST /public/cross-domain/classify}, \texttt{GET /public/cross-domain/results} \\
\hline
Back end (API) & Sign-in, rate limits, storage, reports and safety tools & REST over HTTPS, JSON \\
\hline
Front end & Quick check, analysis, dashboard, tools, cross-domain page & Browser; calls the REST API \\
\hline
Streamlit app & Separate demo of the study & Loads the saved models and never retrains them \\
\hline
Copilot & Answers follow-up questions & REST endpoint calling an outside LLM API, with limits \\
\hline
\end{tabular}
\end{table}

The class diagrams outline the object-oriented structure of the back end. Figure~\ref{fig:class1} shows the application layer: \texttt{MessageService} receives a scan request, stores it through two repositories and calls \texttt{PredictionService}, while \texttt{AuthService} handles sign-in and \texttt{ExtractionService} turns URLs, e-mails, images, PDFs and QR codes into text. Figure~\ref{fig:class2} shows the detection components: \texttt{PredictionService} combines the \texttt{InferenceEngine}, the \texttt{ThreatScorer}, the \texttt{ConfidenceCalculator} and the \texttt{PredictionExplainer}. The engine loads the production model through the \texttt{ModelRegistry} and uses the same preprocessing pipeline and TF-IDF extractor that training used.

\begin{figure}[H]
\centering
\resizebox{\linewidth}{!}{\begin{tikzpicture}[
  cls/.style={draw, rectangle split, rectangle split parts=3, font=\scriptsize, align=left, inner sep=3pt, minimum width=4.2cm, fill=white, rectangle split part fill={gray!18,white,white}},
  rel/.style={-{Stealth[length=2mm]}, thick}]
\node[cls] (auth) at (0,0) {\textbf{\normalsize AuthService}\nodepart{second}+ register(email, password, captcha)\\+ login(...)\\+ refresh\_access\_token(token)\\+ logout(token)\nodepart{third}\textit{issues token pairs}};
\node[cls] (msg) at (6.2,0) {\textbf{\normalsize MessageService}\nodepart{second}- db: Session\\- messages: MessageRepository\\- predictions: PredictionRepository\nodepart{third}+ analyze(user\_id, text, input\_type)\\+ analyze\_ephemeral(text, input\_type)\\+ list\_history(user\_id)};
\node[cls] (mrepo) at (12.2,1.3) {\textbf{\normalsize MessageRepository}\nodepart{second}- db: Session\nodepart{third}+ create(user\_id, text)};
\node[cls] (prepo) at (12.2,-1.9) {\textbf{\normalsize PredictionRepository}\nodepart{second}- db: Session\nodepart{third}+ list\_for\_user(user\_id)\\+ get\_for\_user(id, user\_id)};
\node[cls] (ext) at (0,-3.4) {\textbf{\normalsize ExtractionService}\nodepart{second}URL, e-mail, image, PDF and QR\\inputs turned into text\nodepart{third}\textit{blocks unsafe URLs (SSRF)}};
\node[cls] (ps) at (6.2,-4.0) {\textbf{\normalsize PredictionService}\nodepart{second}- \_engine: InferenceEngine\\- \_confidence\_calculator\\- \_threat\_scorer: ThreatScorer\\- \_explainer: PredictionExplainer\nodepart{third}+ predict(request): PredictionResult};
\draw[rel] (msg) -- node[above, font=\scriptsize] {uses} (mrepo.west);
\draw[rel] (msg.south east) -- node[below, font=\scriptsize, pos=0.55] {uses} (prepo.west);
\draw[rel] (msg.south) -- node[right, font=\scriptsize] {calls} (ps.north);
\draw[rel] (ext.east) -- node[above, font=\scriptsize] {text} (msg.south west);
\end{tikzpicture}
}
\caption{Class diagram of the application layer}
\label{fig:class1}
\end{figure}

\begin{figure}[H]
\centering
\resizebox{0.95\linewidth}{!}{\begin{tikzpicture}[
  cls/.style={draw, rectangle split, rectangle split parts=3, font=\scriptsize, align=left, inner sep=3pt, minimum width=4.2cm, fill=white, rectangle split part fill={gray!18,white,white}},
  rel/.style={-{Stealth[length=2mm]}, thick}]
\node[cls] (ps) at (6.4,0) {\textbf{\normalsize PredictionService}\nodepart{second}- \_engine, \_confidence\_calculator\\- \_threat\_scorer, \_explainer\nodepart{third}+ predict(request): PredictionResult};
\node[cls] (eng) at (0,-3.6) {\textbf{\normalsize InferenceEngine}\nodepart{second}- \_registry: ModelRegistry\\- \_preprocessor, \_vectorizer\\- \_estimator, \_version\_info\nodepart{third}+ load()  + reload()\\+ predict(raw\_text)\\+ transform(raw\_text)};
\node[cls] (ts) at (6.4,-3.6) {\textbf{\normalsize ThreatScorer}\nodepart{second}- \_config: ThreatScoreConfig\nodepart{third}+ intents(raw\_text)\\+ assess(probability, tokens, raw\_text)};
\node[cls] (cc) at (11.0,-3.6) {\textbf{\normalsize ConfidenceCalculator}\nodepart{second}- config: ConfidenceConfig\nodepart{third}+ calculate(scam\_probability)};
\node[cls] (ex) at (12.3,0) {\textbf{\normalsize PredictionExplainer}\nodepart{second}- config: ExplainerConfig\nodepart{third}+ explain(estimator, features,\\\phantom{+ explain(}tokens)};
\node[cls] (reg) at (-0.7,-7.5) {\textbf{\normalsize ModelRegistry}\nodepart{second}versioned model and\\vectorizer files\nodepart{third}+ register()  + promote()\\+ get\_production(name)\\+ list\_versions(name)};
\node[cls] (pp) at (4.6,-7.5) {\textbf{\normalsize TextPreprocessingPipeline}\nodepart{second}cleaner + tokenizer\nodepart{third}applies the same steps in\\training and in serving};
\node[cls] (tf) at (9.6,-7.5) {\textbf{\normalsize TfidfFeatureExtractor}\nodepart{second}word and character n-grams\nodepart{third}+ fit()  + transform(texts)};
\draw[rel] (ps.south west) -- ++(0,-0.5) -| (eng.north);
\draw[rel] (ps.south) -- (ts.north);
\draw[rel] (ps.south east) -- ++(0,-0.5) -| (cc.north);
\draw[rel] (ps.east) -- (ex.west);
\draw[rel] ([xshift=8mm]eng.south west) -- ++(0,-0.4) -| (reg.north);
\draw[rel] (eng.south) -- ++(0,-0.7) -| (pp.north);
\draw[rel] ([xshift=-5mm]eng.south east) -- ++(0,-1.0) -| (tf.north);
\end{tikzpicture}
}
\caption{Class diagram of the detection components}
\label{fig:class2}
\end{figure}

\section{Detailed Design}
\textbf{Use cases.} The use-case diagram in Figure~\ref{fig:usecase} shows how the three kinds of user interact with ScamGuard. A visitor can run a quick check and use the cross-domain page without an account. A registered user can do everything a visitor can and also run a full analysis, read the Scam X-Ray and TrustLens reports, use the safety tools, keep history and cases, and ask the AI copilot, which calls an outside hosted language model. An administrator also sees the registered model versions and can run a robustness test.

\begin{figure}[H]
\centering
\resizebox{0.85\linewidth}{!}{\begin{tikzpicture}[
  uc/.style={draw, ellipse, align=center, font=\footnotesize, minimum width=4.6cm, minimum height=9mm, inner sep=1pt, fill=white},
  act/.style={font=\small\bfseries, align=center},
  ln/.style={thick}]
% system boundary
\draw[thick, rounded corners=3pt] (3.2,-9.3) rectangle (11.4,1.1);
\node[font=\small\bfseries] at (7.3,0.7) {ScamGuard system};
% use cases
\node[uc] (u1) at (7.3,-0.5)  {Quick check (no account)};
\node[uc] (u2) at (7.3,-1.7)  {Cross-domain page: compare\\baseline and adapted models};
\node[uc] (u3) at (7.3,-2.9)  {Register and sign in};
\node[uc] (u4) at (7.3,-4.1)  {Full analysis: text, URL, e-mail,\\screenshot, PDF, QR code};
\node[uc] (u5) at (7.3,-5.3)  {View Scam X-Ray and TrustLens};
\node[uc] (u6) at (7.3,-6.5)  {Use safety tools, history and cases};
\node[uc] (u7) at (7.3,-7.7)  {Ask the AI copilot};
\node[uc] (u8) at (7.3,-8.8)  {View model versions, run robustness test};
% actors (stick figures)
\foreach \n/\y/\l in {v/-1.2/Visitor, r/-4.6/Registered\\user, a/-8.0/Administrator}{
  \draw[thick] (0.9,\y+0.55) circle (0.2);
  \draw[thick] (0.9,\y+0.35) -- (0.9,\y-0.2);
  \draw[thick] (0.5,\y+0.15) -- (1.3,\y+0.15);
  \draw[thick] (0.9,\y-0.2) -- (0.55,\y-0.65);
  \draw[thick] (0.9,\y-0.2) -- (1.25,\y-0.65);
  \node[act] at (0.9,\y-1.1) {\l};
}
% external actor
\draw[thick] (13.1,-6.65) circle (0.2);
\draw[thick] (13.1,-6.85) -- (13.1,-7.4);
\draw[thick] (12.7,-7.05) -- (13.5,-7.05);
\draw[thick] (13.1,-7.4) -- (12.75,-7.85);
\draw[thick] (13.1,-7.4) -- (13.45,-7.85);
\node[act] at (13.1,-8.35) {Hosted LLM\\(Gemini)};
% associations
\draw[ln] (1.5,-1.2) -- (u1.west); \draw[ln] (1.5,-1.2) -- (u2.west); \draw[ln] (1.5,-1.2) -- (u3.west);
\draw[ln] (1.5,-4.6) -- (u4.west); \draw[ln] (1.5,-4.6) -- (u5.west); \draw[ln] (1.5,-4.6) -- (u6.west); \draw[ln] (1.5,-4.6) -- (u7.west);
\draw[ln] (1.5,-8.0) -- (u8.west);
% generalisation arrows between actors
\draw[thick, -{Triangle[open, length=3mm]}] (0.9,-3.2) -- (0.9,-2.45);
\draw[thick, -{Triangle[open, length=3mm]}] (0.9,-6.6) -- (0.9,-5.95);
\draw[ln, dashed] (u7.east) -- (12.55,-7.0);
\end{tikzpicture}
}
\caption{Use-case diagram of ScamGuard}
\label{fig:usecase}
\end{figure}

\textbf{Request flow.} The browser and the back end exchange JSON over HTTPS. The sequence diagram in Figure~\ref{fig:seq} traces a single analysis from request to reply.

\begin{figure}[H]
\centering
\resizebox{0.95\linewidth}{!}{\begin{tikzpicture}[font=\small, >={Stealth[length=2.2mm]}, msg/.style={->, thick}, ret/.style={->, thick, dashed}, lbl/.style={above=1pt, fill=white, inner sep=1.5pt, font=\small}]
\foreach \x/\n/\t in {0/u/User, 3.8/c/Web client\\(Next.js), 7.8/a/FastAPI\\back end, 11.8/m/Detector, 15.6/d/Database} {
  \node[draw, rounded corners=2pt, fill=gray!12, minimum width=2.8cm, minimum height=12mm, inner sep=3pt, align=center] (\n) at (\x,0) {\t};
  \draw[gray] (\x,-0.6) -- (\x,-11.0);
}
\draw[msg] (0,-1.5) -- node[lbl]{1. paste message} (3.8,-1.5);
\draw[msg] (3.8,-2.6) -- node[lbl]{2. POST /analyze} (7.8,-2.6);
\node[draw, fill=white, align=center, font=\footnotesize, inner sep=3pt, minimum width=3.2cm] at (7.8,-3.8) {3. check token, rate\\and size limits; clean};
\draw[msg] (7.8,-5.1) -- node[lbl]{4. predict(text)} (11.8,-5.1);
\node[draw, fill=white, align=center, font=\footnotesize, inner sep=3pt, minimum width=3.0cm] at (11.8,-6.3) {5. classifier, rules,\\fusion, explanation};
\draw[ret] (11.8,-7.6) -- node[lbl]{6. verdict, score, phrases} (7.8,-7.6);
\draw[msg] (7.8,-8.4) -- node[lbl]{7. save scan} (15.6,-8.4);
\draw[ret] (15.6,-9.1) -- node[lbl]{8. saved} (7.8,-9.1);
\draw[ret] (7.8,-9.9) -- node[lbl]{9. 200 OK (JSON)} (3.8,-9.9);
\draw[ret] (3.8,-10.7) -- node[lbl]{10. show verdict} (0,-10.7);
\end{tikzpicture}
}
\caption{Sequence diagram of one message analysis}
\label{fig:seq}
\end{figure}

\textbf{Interfaces.} The FastAPI back end exposes REST routes for sign-in, scans, history, cases, analytics, the copilot, the safety tools and the cross-domain study. The Next.js front end calls them through its own origin and offers the quick check, full analysis, dashboard, safety tools and cross-domain page.
\begin{itemize}
\item \textbf{Endpoints.} The three classification routes are listed in Table~\ref{tab:endpoints}. Further routes cover sign-in, history, cases, analytics, the copilot and the safety tools.
\item \textbf{Authentication.} Short-lived signed access tokens paired with refresh tokens; passwords kept as bcrypt hashes.
\item \textbf{Input limits.} A message may be up to 4,000 characters long, or 2,000 on the public endpoints.
\end{itemize}

\begin{table}[H]
\centering
\caption{Classification endpoints and their rate limits}
\label{tab:endpoints}
\small
\begin{tabular}{|l|l|c|}
\hline
\textbf{Route (POST)} & \textbf{Used by} & \textbf{Requests/min} \\
\hline
\texttt{/api/v1/public/quick-check} & Visitors (no account) & 10 \\
\hline
\texttt{/api/v1/messages/analyze} & Signed-in users & 30 \\
\hline
\texttt{/api/v1/public/cross-domain/classify} & Cross-domain page & 30 \\
\hline
\end{tabular}
\end{table}

\section{Data Design}\label{sec:datadesign}
\textbf{Database.} The main tables of the PostgreSQL database are listed in Table~\ref{tab:db}. Every query is scoped to the signed-in user.

\begin{table}[H]
\centering
\caption{Main database tables}
\label{tab:db}
\begin{tabular}{|L{4.6cm}|L{9.7cm}|}
\hline
\textbf{Table} & \textbf{Contents} \\
\hline
\texttt{users} & Account, hashed password, role, preferences \\
\hline
\texttt{refresh\_tokens} & Issued refresh tokens for rotation and revocation \\
\hline
\texttt{messages} & Each analysed input with its type (text, URL, e-mail, image, PDF, QR) \\
\hline
\texttt{predictions} & Verdict, classifier probability, risk score, category, explanation, model version \\
\hline
\texttt{cases}, \texttt{case\_scan\_links}, \texttt{case\_timeline\_entries} & Investigation cases, the scans linked to them and their notes \\
\hline
\texttt{community\_reports} & One report per user for a number, link domain, UPI ID or e-mail \\
\hline
\texttt{alerts}, \texttt{audit\_log}, \texttt{model\_registry\_meta} & Alerts, security audit trail and registered model versions \\
\hline
\end{tabular}
\end{table}

\textbf{Dataset files.} Each row of \textbf{ScamGuard-IN-2026} has seven fields: \texttt{text}; \texttt{label}, which is scam or legitimate; \texttt{source}; \texttt{subtype}; \texttt{is\_synthetic}; \texttt{campaign}, naming the near-duplicate group; and \texttt{split}, which is train, development, test or 2026 test.

After cleaning, the \textbf{cross-domain datasets} keep three fields: \texttt{message}, \texttt{label} (spam = 1, ham = 0) and \texttt{domain} (UCI or India). Alongside them the study saves every result table as a CSV file, the error sheet with its hand-filled \texttt{manual\_category} column, a \texttt{results.json} summary and four joblib models, namely Naive Bayes and Logistic Regression, each as a baseline and as an adapted version.

\textbf{Model registry.} Each trained model is stored as a versioned pair of files, the fitted TF-IDF vectorizer and the classifier, with an index entry that records its version and metrics. The back end loads the production version once at start-up.
